% Generated draft: requires mathematical and source review.
\documentclass[11pt]{article}
\usepackage[margin=25mm]{geometry}
\usepackage{fontspec}
\setmainfont{Latin Modern Roman}
\usepackage{amsmath,amssymb,mathtools,unicode-math}
\setmathfont{Latin Modern Math}
\usepackage{enumitem,needspace,xurl,hyperref}
\setlength{\parindent}{0pt}
\setlength{\parskip}{5pt}
\setlength{\emergencystretch}{3em}
\newcommand{\N}{\mathbb{N}}
\newcommand{\Z}{\mathbb{Z}}
\newcommand{\Q}{\mathbb{Q}}
\newcommand{\R}{\mathbb{R}}
\newcommand{\C}{\mathbb{C}}
\newcommand{\F}{\mathbb{F}}
\newcommand{\A}{\mathbb{A}}
\newcommand{\PP}{\mathbb{P}}
\newcommand{\E}{\mathbb{E}}
\newcommand{\eps}{\varepsilon}
\newcommand{\dd}{\,\mathrm{d}}
\begin{document}
\section*{UnsolvedMath ID: 30005297}
\textbf{Source problem code:} catalogue entry
\subsection{Definitions and mathematical statement}
The supplied excerpt does not specify the ambient space, ground field for the point configurations, or meaning of ``nontrivial.'' The following gives a conventional formulation in complex projective three-space, conditional on that being the intended setting; the omitted restrictions require mathematical review.

Let \(Z\subset\PP^3(\C)\) be a finite set of distinct points, regarded as a reduced scheme. For positive integers \(a,b\), call \(Z\) \((a,b)\)-geproci if there is a nonempty Zariski open set of projection centers \(o\in\PP^3(\C)\setminus Z\) such that projection \(\pi_o\) to a projective plane is injective on \(Z\) and its image is a scheme-theoretic complete intersection
\[
\pi_o(Z)=V(F_o,G_o)\subset\PP^2(\C),\qquad \deg F_o=a,\quad\deg G_o=b,
\]
where the two homogeneous polynomials have no common curve component. Call \(Z\) geproci if such \(a,b\) exist. Necessarily \(|Z|=ab\). This definition permits all positive degrees; the source does not specify which degrees or configurations its word ``nontrivial'' excludes.

The underlying matroid \(M(Z)\) has ground set \(Z\) and rank function
\[
r_Z(S)=\dim_{\C}\operatorname{span}\{v_p:p\in S\},
\]
where \(v_p\in\C^4\setminus\{0\}\) represents \(p\). This is independent of the choices of representatives. Two configurations realize the same matroid if a bijection between their points preserves the rank of every subset. The uniform matroid \(U_{r,n}\) has rank function \(S\mapsto\min\{|S|,r\}\) on an \(n\)-element ground set. For a spanning configuration in \(\PP^3(\C)\), linear general position means that its matroid is \(U_{4,n}\).

A finite complex reflection group is a finite subgroup \(W\leq\operatorname{GL}(V)\), for a finite-dimensional complex vector space \(V\), generated by nonidentity elements fixing a hyperplane pointwise. It is irreducible if \(V\) has no nonzero proper \(W\)-invariant subspace.

Subject to the unresolved scope qualifications above, the original excerpt poses these questions:
\begin{enumerate}
\item Can the classification of irreducible complex reflection groups provide ideas toward a classification of nontrivial geproci sets? The excerpt specifies neither a construction associating point sets to reflection groups nor a proposed classification theorem.
\item What necessary conditions on \(M(Z)\), beyond elementary requirements such as a suitable factorization \(|Z|=ab\), must hold when \(Z\) is geproci?
\item Does there exist a geproci realization of a uniform matroid, described in the source as a set of points in linear general position? In the spanning three-dimensional interpretation, this asks for a configuration with \(M(Z)\cong U_{4,n}\); the intended restrictions on \(n,a,b\) and on trivial examples are unspecified.
\item Is the geproci property determined by the underlying matroid? For configurations in the intended class, is
\[
M(Z)\cong M(Z')\quad\Longrightarrow\quad
\bigl(Z\text{ is geproci}\ \Longleftrightarrow\ Z'\text{ is geproci}\bigr)
\]
always true? Here ``geproci'' is existential in the degrees; the source does not ask explicitly for preservation of a fixed pair \((a,b)\).
\end{enumerate}
The supplied shortened statement additionally asks which matroid properties are invariant under ``the resulting reflection constructions.'' No such constructions are defined in the record, so this additional wording cannot yet be formulated as a precise invariance assertion.

The source compares the matroid question with Terao's freeness conjecture. It says that a direct connection between geproci and freeness is not expected, while suggesting projective duality, Matlis duality, or other algebraic considerations as possible ways for arrangement theory to inform the problem. These are motivations, not asserted implications.
\subsection{Short English statement}
The source asks whether complex reflection groups can guide a classification of nontrivial geproci point sets, which combinatorial restrictions geproci sets satisfy, whether uniform matroids admit geproci realizations, and whether the geproci property depends only on the underlying matroid. A conventional formulation in \(\PP^3(\C)\) is given, but the intended ambient setting, exclusions implicit in ``nontrivial,'' and reflection constructions remain unspecified.
\subsection{Sources}
\begin{itemize}
\item Supplied local UnsolvedMath registry record: permanent ID \(30005297\), source problem code \texttt{catalogue entry}, title ``Matroid Invariance and Reflection Constructions for Geproci Sets,'' in the Oberwolfach Reports Open Problems collection. No individual proposer is supplied; the proposed year is recorded as 2022.
\item Supplied source citation: \emph{Mini-Workshop: Subvarieties in Projective Spaces and Their Projections} (2023), Oberwolfach Reports, pp.\ 3093--3119, DOI \texttt{10.4171/owr/2022/54}. The excerpt carries page number 3095. \url{https://doi.org/10.4171/owr/2022/54}. The citation's year 2023 and the record's proposed year 2022 are retained as supplied.
\item Supplied literature source: Chiantini et al., \emph{Geproci sets and skew-line combinatorics}. \url{https://arxiv.org/abs/2308.00761}.
\item Supplied literature source: EMS Press, OWR 2022/54. \url{https://doi.org/10.4171/owr/2022/54}.
\item Supplied dated literature curation, checked at \texttt{2026-08-22}: the record labels the reflection-group, uniform-matroid, and same-matroid questions open. It reports that known configurations realizing the same matroid need not be projectively equivalent, which does not settle geproci invariance. This assessment is supplied dated curation, not independently verified; the background expressly requests independent verification of time-sensitive status and cited claims.
\item The registry labels the statement \texttt{corrected\_verified} and reports comparison with the official Oberwolfach report. This is a supplied verification claim, not an independent verification performed for this draft; the difference between the original excerpt and shortened statement remains recorded below.
\end{itemize}
\end{document}
